% !TeX root = ../教材题库_全部条目_转写工作稿.tex
% Source: XJ-B1-C04-PR-P097-Y1020
% 第4章 / 4.1 实数指数幂和幂函数 / 教材 p.97 / PDF 第104页
% OCR未校对：公式与文字均可能有误；未完成数学TeX转写；禁止直接发给学生
% 原题及相邻原解截图（核对用）：../题目原图/XJ-B1-C04-PR-P097-Y1020-01.webp
% 本文件只在转写完成后作为题干使用；OCR 条目尚不是可靠题干。
\begin{bankitem}{p.97 练习 第2题}{第4章 / 4.1 实数指数幂和幂函数 / 教材 p.97 / PDF 第104页}{XJ-B1-C04-PR-P097-Y1020}{OCR未校对：公式与文字均可能有误；未完成数学TeX转写；禁止直接发给学生}
\begin{ocrdraft}
[教材p.97 / PDF第104页]
\par
2．用分数指数幂的形式表示下列各式（式中字母都是正数）：
\par
（2）
\par
（m十7D3
\par
（3）
\par
十2b）2\ensuremath{\cdot}
\par
（4）
\par
（5）
\par
a乃．
\par
（6）
\end{ocrdraft}
\end{bankitem}
